% !TeX spellcheck = en_US
\documentclass[french]{yLectureNote}

\title{Algèbre linéaire 3*}
\subtitle{Physique}
\author{Paulhenry Saux}
\date{\today}
\yLanguage{Français}
%lombardi@math.univ-toulouse.fr
\professor{E.Lombardie}
\usepackage{graphicx}%----pour mettre des images
\usepackage[utf8]{inputenc}%---encodage
\usepackage{geometry}%---pour modifier les tailles et mettre a4paper
%\usepackage{awesomebox}%---pour les boites d'exercices, de pbq et de croquis ---d\'esactiv\'e pour les TP de PC
\usepackage{tikz}%---pour deiffner + d\'ependance de chemfig
% \usepackage{tabularx}%---pour dimensionner automatiquement les tableaux avec variable X
\usepackage{awesomebox}%---Pour les boites info, danger et autres
\usepackage{menukeys}%---Pour deiffner les touches de Calculatrice
\usepackage{fancyhdr}%---pour les en-t\^ete personnalis\'ees
\usepackage{blindtext}%---pour les liens
\usepackage{hyperref}%---pour les liens (\`a mettre en dernier)
\usepackage{caption}%---pour la francisation de la l\'egende table vers Tableau
\usepackage{pifont}
\usepackage{array}%---pour les tableaux
\usepackage{lipsum}
\usepackage{yFlatTable}
\usepackage{multicol}
\newcommand{\Lim}[1]{\lim\limits_{\substack{#1}}\:}
\renewcommand{\vec}{\overrightarrow}
\newcommand{\N}[0]{\mathbb{N}}
\newcommand{\R}[0]{\mathbb{R}}
\newcommand{\dd}{\mathrm{d}}
\newcommand{\norm}[1]{||\vec{#1}||}
\newcommand{\fo}{\psi(\vec{r},t)}
\newcommand{\foe}{\psi(\vec{r},t)\*}
\newcommand{\HH}{\hat{H}}
\newcommand{\hb}{\hbar}
\newcommand{\lap}{\nabla^2}
\newcommand{\lapcc}{\frac{\partial^2 }{\partial x^2}+\frac{\partial^2 }{\partial y^2}+\frac{\partial^2 }{\partial z^2}}
\newcommand{\E}{\vec{E}(M)}
\newcommand{\B}{\vec{B}(M)}
\newcommand{\V}{V(M)}
\newcommand{\er}{\vec{e_{\rho}}}
\newcommand{\ep}{\vec{e_{\varphi}}}
\newcommand{\pip}{\pi^+}
\newcommand{\pim}{\pi^-}
\newcommand{\mv}{\mathcal{V}}
\DeclareMathOperator{\Div}{Div}
\newcommand{\rot}{\vec{\mathrm{rot}}}
\newcommand{\grad}{\vec{\mathrm{grad}}}
\begin{document}
\setcounter{chapter}{1}
%C2 : Applications linéaires sur espaces euclidiens
%C3 : Espaces hermitiens
%C4 : Formes quadratiques
%C5 : Introduction aux groupes

%TD = U-111/U-112
%CM : B-08/B-07/Grignard

%CC2 : TD du 19 mars
%CC1 : 1h30
%CC2 : TD du 7 mai
%CC3 : 2 juin
%CC4 : 9 juin
\chapter{Endomorphismes d'espaces euclidiens}
$S(x) = x_f-x_g$ avec F et G 2 espaces supplémentaires. On décompose $x = x_f-x_g$ et $P_{F//G}=x_f$.

Notions sur les symétries à revoir
\section{Isométrie et matrices orthogonales}
\subsection{Définitions et propriétés}
\begin{definition}[Isométrie]
    $f : E\to E$ est une isométrie quand $\forall x\in E, ||f(x)|| = ||x||$
\end{definition}
\begin{proposition}
    $\forall x \in E, ||f(x)||  = ||x|| \iff (f(x)|f(y)) = (x|y) \iff $ pour toites BON, f des vecteurs est
    aussi BON
\end{proposition}
\begin{myproof}
    $1\rightarrow 2$ :

    $(x|y) = \frac12 (||x+y||^2-||x||^2-||y||^2) = \frac12 (||f(x+)y||^2-||f(x)||^2-||f(y)||^2)$
    car f conserve la norme par hypothèse.

    Mais f est linéaire, donc $(x|y) =\frac12 (||f(x)+f(y)||^2-||f(x)||^2-||f(y)||^2)$
    Donc $1\rightarrow 2$ 

    $2\rightarrow 3$ 

    On supposera f conserve le froduit scalaire. Soit $(e_1,\dots, e_n)$ une BON.
    Donc $\forall i,j (e_i|e_j)=\delta_{ij}\iff \forall i,j (f(e_i)|f(e_j))=\delta_{ij}\iff $ BON
    Donc $2\rightarrow 3$ 

    $3\rightarrow 1$ 

    f envoie une BON et on veut montrer que f conserve la norme.

    On écrit $x=x_1e_1+\dots+x_ne_n$. De plus, comme $f$ est linéaire, $f(x) = x_1f(e_1)+\dots+x_nf(e_n)$

    $||x^2|| =  x_1^2+\dots+x_n^2$ car BON.
    
    $|f(x)|| =  x_1^2+\dots+x_n^2$ car BON.
\end{myproof}
\begin{definition}
    Une matrice $Q\in M_n(\R)$ est orthogonale ssi son inverse $Q^{-1} = ^TQ$
\end{definition}
Exemples : Rotation
\begin{proposition}
    E euclidien. f est une isométrie $\iff Mat(f)_e$ est orthogonale.
\end{proposition}
\criticalInfo{Condition}{C'est faux si ce n'est pas une BON}
\begin{myproof}
    f est isométrie .

    A = Mat(f). On décompose x et y dans la base.

    $(x|y) = x_1y_1+\dots+x_ny_n = (^tX)Y$ car c'est BON et on utilise des vecteurs colonne

    $f(x)|f(y) = (x|y) \rightarrow ^t(AX)(AY) = ^tXY \iff ^tX^tAAY = ^tXY \iff  ^tX(^tAAY-Y) = 0$ 
    
    $\iff (X|^tAAY-Y) = 0 $
    $\iff ^tAAY-Y \in (\R^n)^{\perp}$

    $ \iff ^tAAY-Y = 0 \iff ^tAA-I_n = I_n \iff A^{-1}  = ^tA$


\end{myproof}
\begin{proposition}
    Si Q matrice orthogonale, alors $\det(Q) = \pm 1$
\end{proposition}
\begin{myproof}
    $^tQQ = I_n$, donc $\det(^tQQ) = \det(I_n) = 1 = \det(Q)\det(^tQ)$ et comme $\det(^tQ) = \det(Q)$, c'est bien 1 ou -1.
\end{myproof}
$det(\lambda A) = \lambda^n A$.
\begin{proposition}
    Si f est une isométrie, alors les valeurs propres réelles de f sont $\pm 1$ $\iff$ Si A est orthogonale alors les vp réelles
    de A sont $\pm 1$
\end{proposition}
\begin{myproof}
    Si $\lambda \in \R,$ vp de f, $\exists x\neq 0, f(x)=\lambda x$

    Donc $||f(x)|| = ||\lambda x||$ donc $|\lambda| = 1$.
\end{myproof}
\begin{definition}
    On note $O(n)$ l'ensemble des matrices orthogonales et $SO(n)$ l'ensemble des matrices
    orthogonales Q telle que $\det(Q) = 1$
\end{definition}
\begin{proposition}
    Si (A,B) appartiennent à O(n), leur produit aussi. In appartient à O(n).
    Si A appartient à O(n), l'inverse aussi.
\end{proposition}
Ces trois propriétés assurent que O(n) est un groupe multiplicatif.
\begin{proposition}
    SO(n) est aussi un groupe multiplicatif
\end{proposition}
On dit que $O(n)$ est le groupe orthogonal et SO(n) le groupe spécial orthogonal.
\begin{myproof}
    Pour le produit, on fait un calcul de AB et on obtient bien l'identité

    Pour la 3eme, avec règle de transposé, etc pour montrer que $^t(A^{-1}) = (A^{-1})^{-1}$
\end{myproof}
\subsection{Étude de O(2)}
\begin{proposition}
    Soit A dans $O(2)$. Alors soit $A\in SO(2)$ et dans ce cas A = (cos -sin//sin cos)
    A est une rotation d'angle theta
    Soit A n'appartient pas à SO(2) et on a $A = (cos sin // sin -cos)$ A est une symétrie orthogonale d'angle theta/2
\end{proposition}
\begin{myproof}
    Pour le premier cas, on fait en posant A = (a b // c d )

    On obtient comme équations $a^2+c^2=1, b^2+d^2=1, ab+cd=0$
    
    Ils ont la forme d'un cercle : $a=\cos \theta, c=\sin\theta, b=\cos\phi, d = \sin \phi, \cos(\phi-\theta)=0$

\end{myproof}
\subsection{Étude de O(3)}
\begin{proposition}
    A appartient à O(3) et f l'endomorphisme de $\R^3$ tel que $A = mat(f,e_i)$ avec $e_i$ la base canonique de $\R^3$.

    Alors il existe une bon de $\R^3$ telle que  $A' = (cos -sin 0//sin cos 0//0 0 epislon)$
    avec epsilon =1 si A appartient à SO(3) et -1 sinon.
\end{proposition}
\begin{myproof}
    Étape 1 : Lemme : Si F est une isométrie et si F un sev de E stable par $f$. Alors :
    alors $f(F^{\perp})\subset F^{\perp}$

    Cela fait instananément des matrices diagonales par bloc

    Exercice pour démontrer le lemme

    Étape 2 : On veut trouver un sous espace stable à l'aide d'un vecteur propre réel.
    $\lambda v_p \iff P_A(\lambda) = 0$

    Le dégré vaut 3, donc P a au moins une racine réelle.

    finir la demo
\end{myproof}
\begin{proposition}
    Soit A appartient à O3 et différent de l'identité.
    a) Soit det A = 1, A représente une rotation autour de l'axe $E_1$ est donné par $tr A = 2\cos(\theta)+1$
    
    b) Si det A = -1, A représente une rotation d'axe $E_{-1}$ d'angle théta non orieté donné par $tr A = 2\cos theta -1$
\end{proposition}
\subsection{Classification dans O(n)}
\begin{theorem}
    Soit f isométrie de E espace euclidien. Alors il existe une bon de E tq A  = mat(f,e) est de la forme :

    avec r_theta = [cos -sin // sin cos]
\end{theorem}
\begin{theorem}
    Soit E euclicien, f isométrie de E. On note $E_1 = \{x\in E/ f(x) = x\}$

    Alors f s'écrit $f = \sigma,\dots \sigma_\theta$
    n- dim E_1 < \sigma < n
    \sigma réflexion cad une symétrie orthogonale par rapport à un hyperplan
\end{theorem}
H hyperplan de E : $H sev de E, dim H  = n-1$

Dans une bon adapatée à la matrice e l'une est de la forme : S = [1--0-0/0-0---1//-1]
%----------------------Cours précédent non travaillé ---------------
\section{Adjoint d'un endomorphisme}
\subsubsection{Endomorphisme adjoint}
\begin{definition}
    Soit E un espace euclidien et f une application linéaire. Il existe un unique endomorphisme de E, noté f* tel que 
    $$\forall (x,y)\in E, (f(x)|y) = (x|f*(y))$$
    $f*$ est appelé adjoint de f.

    Si $(e_i)$ est une BON, alors $A = mat(f,e), A*=mat(f*,e), A* = ^tA$
\end{definition}
\begin{myproof}
    $\forall x,y \in E (f(x)|y)=(x|f*(y)) \iff ^t(AX)Y=^tX(A*Y) \iff (^tA-A*)Y=0$
\end{myproof}
\begin{proposition}
    \begin{itemize}
        \item $f** = f$
        \item Id* = Id
        \item $\lambda f* = (\lambda f)*$
        \item $(f\circ g)* = g*\ciric f*$
        \item Si f inversible : $(f*)^{-1} = (f^-1)$
    \end{itemize}
\end{proposition}
\begin{proposition}
    \begin{itemize}
        \item $rg f* = rg f$
        \item $det f* = det f$
        \item tr f* = tr f
    \end{itemize}
\end{proposition}
\begin{proposition}
    Soit E euclidien et $F\subset E$, f app lin. On a 
    
    Si F est stable par f $\iff f^{\perp}$ stable par $f*$.
\end{proposition}
\begin{myproof}
    Étape 1 : $\rightarrow$

    Soit $y\in F^{\perp}$ on veut montrer que $f*(y)\in F^{\perp}$

    $\forall x\in  F, (x|f*(y)) = (f(x)|y) = 0$

    Donc $f*(y)\in F^{\perp}$, donc $f*(F^{\perp})\subset F^{\perp}$ donc $F^{\perp}$ est stable par $f*$.

    On fait l'autre sens :
    $F^{\perp}$ stable par $f*$ donc $(F^\perp)^\perp$ stable par $f**$ donc F stable par $f$.

\end{myproof}
\begin{proposition}
    $Ker f* =( Im f)^{\perp}$

    $Im f* = (Ker f)^{\perp}$
\end{proposition}
\begin{myproof}
    $y\in Ker f* \iff f*(y) = 0\iff f*(y)\in E^{\perp} \iff (x|f*(y)) = 0 \iff (f(x)|y) = 0\iff y\in (Im f)^{\perp}$
\end{myproof}
\section{Endomorphisme Auto-adjoint}
\begin{definition}
    f est dit auto-adjoint $\iff \forall x,y, \in E , (f(x)|y) = (x|f(y))\iff f=f*$
\end{definition}
Pour une matrice, c'est l'ensemble des matrices symétriques
\begin{proposition}
    E espace euclidien, A = mat(f,e(BON)), f auto-adjoint $\iff A$ symétrique.
\end{proposition}
\begin{theorem}
    Si f est auto-adjoint alors

    1. f est diagonalisable et toutes les valeurs propres sont réelles
    2. Les sous-espaces propres sont deux à deux orthogonaux. On peut donc construire une BON de E, constituée de vecteurs propres de f en prenant une BON sur chaque espace propre

\end{theorem}
\begin{theorem}
    SI A est symétrique, alors A diagonalisable en BON : $\exists Q \in O_n tq A'=Q^{-1}AQ = ^tQAQ$ est diagonale.
\end{theorem}
\checkInfo{Rappel}{A diagonalisable si il existe P tq $A'=P^{-1}AP$ est diagonale. }
\begin{myproof}
    Étape 1 :Montrer que toutes kes vp sont réelles.

    Soit f auto-adjoint, e BON, A=Mat(f,e), tA = Mat(f,E)

    $\lambda v_p$ de f. c'est une racine du polynome caractétistique de A. On suppose que lambda est complexe.

    On multiplie par le $X bar$ qui est le produit scalaire complexe.

    Donc $\lambda = \bar{\lambda}$

    On montre que c'est diagonalisable.

Étape 2 : Montrons que f est diagonalisable par recurrence.

Initialisation : n=1

Supposons la propriété vraie au rang n.  Soit E avce dim E = n+1 et f est une app lin auto-adjointe.

Donc $\exists \lambda vp$ Soit $e_{n+1}\neq 0$ et $f(e_n) = \lambda e_{n+1}$

F = R e_{n+1} et f(F) \seubset F. Donc D^\perp stable par f* = f

En considérant $f_G : G\to G$ on constante que $d_m G = n, f_G $ est auto-adjointe.

Donc il existe e_n base de G constituée.

donc il existe une base de G contrituée de cecteurs propres de f. Don c'est une base de E.

Étape 3 : Montrons que les sous espaces propres sont orthogonaux entre eux.

Étape 3 à savoir :
Soit $\lambda_2,\lambda_1$ deux vp diffétentes : $f(x_1)=\lambda_1x_1$.

$\lambda_1(x_1|x_2) = (f(x_1)|x_2) = (x_1|f*(x_2)) = (x_1|f(x_2)) = \lambda_2(x_1|x_2) \Rightarrow (\lambda_1-\lambda_2)(x_1|x_2) = 0$ et les lambdas sont différentes donc $(x_1|x_2)=0$
\end{myproof}
\checkInfo{Méthode pour diagonalisr en BON une matrice symétrique}{
    \begin{itemize}
        \item Véirier que la matrice est symétrique
        \item Calcul du polynome caractétistique
        \item Calcul des racines (= valeur propre et multiplicité)
        \item Pour chaque kalbdan trouver une base quelconque
        \item Pour chaque base utiliser GS
    \end{itemize}
}
%----------------- Cours du 24/03/2026---------------

\subsection{Endomorphisme auto-adjoint positif et décomposition polaire}
\begin{definition}
    Soit E euclidien, et $f$ auto-adjoint. $f$ est positive $\iff (x|f(x))\geq 0$ Elle est définie positive $\iff x\neq 0 , (x|f(x))>0$.
\end{definition}
\begin{proposition}
    $f$ est auto-adjoint. $f$ spotive $\iff$ toutes les valeurs propres sont positives. Elle est déifnie positive $\iff$ toutes les valeurs propres sont strcutement positives.
\end{proposition}
\begin{myproof}
    $f$ auto-adjoint, donc diagonalisable en base orthonormée de valeurs propres réelles. Donc in y a une BON $e_i$ et $f(e_i) = \lambda_i e_i, \lambda_i \in \R$

    on  a $x= \sum x_i e_i$ , $f(x_i) = \sum \lambda_i x_i e_i$, donc $(x|f(x)) = \sum (x_i)(\lambda_i)(x_i) = \sum \lambda_i x_i^2$

    a) $f$ positif $\iff \forall(x_i) \iff (\lambda_1 x_1+\dots+\lambda_n x_n\geq 0)\iff \lambda_i \geq 0$
\end{myproof}
\begin{theorem}
    Soit $M\in Gl_n(R)$. Alors il existe un unique couple $(Q,S)$ où Q est une matrice orthogonale et S symétrique définie positive tel que $M=QS$.
\end{theorem}
\begin{lemma}
    Soit A symétrique  définie positive. Alors il existe une unique matrice S symétrique définie positive tel que $S^2=A$.
\end{lemma}
\begin{myproof}[Théorème]
    Étape 1 : Montrons que $^tMM$ est symétrique définie positive.
On a $^tX(^tMM)X = ^tYY$ avec $Y = MX$

De plus, $^tX(^tMM)X \geq 0$
et $^tX(^tMM)X = 0 \iff Y = 0 \iff MX = 0 \iff X = 0$

Étape 2 : Analyse :
Si $M=QS$, alors $^tM = ^tS^tQ = S^tQ$

Alors $S = \sqrt{^tMM}, Q = MS^{-1}$

Étape 3
M = QS, OK, S est définie positive, Ok 

Reste à montrer que $Q \in O(n)$..

$^tQQ = ^t(MS^{-1})MS^{-1} = I_n$.
\end{myproof}
\begin{myproof}[Démonstration du lemme]
    Existence de $S^2 =A$

    On peut diagonaliser A.

    $\exists Q_1 \in O(n), ^tQ_1Aq_1 = D$

    On peut prendre $S = Q_1 \sqrt{D}^tQ_1$ qui convient (A dem)

    Montrons maintenant l'unicité :

    Montrons que si $S^=A$, alors $\exists e_i$ bon de vecteurs propre de A et S.

    point clef pour faire la co-réduction de A et B, c'est que A et B commutent.
    Donc si $S^2=A, SA=AS$

    A symétrique donc soit ei une BON de vecteurs propres de A.

    $A S e_i = SA ei \iff A S_ei = \lambda_i S e_i$.

    En notant $E_{\lambda}$ le sous espace propre de A associé à $\lambda_i$, on a $S(E\lambda_i )\subset E\lambda_i$

    Sur chaque $E\lambda_i, S$ est symétrique définie positive donc diagonalisable en BON. En prenant sur chaque $E\lambda_i$ une base qui diagonalise $S^{-1}$ et concaténant ses bases on obtient une BON de vecteurs propres de  et de S.

    En notant $P$ la matrice de passage, on a $S^2=A\iff ^tpS^2P = ^tPAP$

    Il y a une unique solution donnée par $\lambda_i ' = \pm \sqrt{\lambda_i}$.

\end{myproof}
\begin{proposition}[Corollaire]
    Si $M$, alors il existe $Q\in o(n)$ et $S$ symétrique positive tq $M=QS$ et Q et S ne sont pas forcément unique.
\end{proposition}
\subsection{Valeurs singulières}
Dans cette section, on travaille dans Mn(R), matrices rectangles de taille $n\times m$ à coefficients réels. n = nombre de ligne (nombre de vecteurs de l'espace d'arrivé) et m nombre de colonne (nombre de vecteurs de l'espace de départ)
\begin{lemma}
    Soit M\in Mn,m(R)
    a) $^tMM^$ est une matrice symétrique positive de taille m*m.
    b) $Ker(^tMM) = Ker (M)$
    c) rg(^tMM) = rg(M)
    d) ^tMM = 0 \iff M = 0
\end{lemma}
\begin{myproof}
    a) ^tMM est bien de taille m*m symétrique et $^tX^tMMX = ^t(MX)MX = (x|^tMMX) = ||MX||\geq 0$
    b) Si X\jn ker(M), MX=0 \rightarrow ^tMMX = 0, donc \lambda \in ker^tMM, donc ker M\subset ker(^tMM)
    On peut faire la réciproque.
    c) rg(M) = dim(Im(M)) = n - dim(ker M) = rg(^tMM)
    d) M = 0 \iff ker(M) = R^m \iff ^tMM = 0
\end{myproof}
\begin{theorem}[Décompoisiton en valeur singulières]
    Soit A\in M_{n,m}(R), r= rg(A). Alors $Q\in O(n), P\in O(n)$ et $G\sigma_1,\dots, \sigma_r >0$ tel que $A = Q\sum ^tp$ où $\sum$ matrices $n\times m$ de ma forme et \sum = [\sigma_1 --- sigma_r] en diagonale. Ici la somme n'est pas une somme mais une matrice.
\end{theorem}
\begin{myproof}
    /!\ démonstration inutile
    content\begin{enumerate}
        \item ^tAA \in M_m R symétrique positive  donc il existe P\in O(m) tq ^tp^tAAP  = [\lmabda_10//0\lamnda_m] où \lambda_1 >0. On peut imposer \lambda_1\geq \lambda_2\geq \dots\geq \lambda_n \geq 0
        Donc \lambda_1 \geq \lambda_2\geq\dots\geq \lambda_R
        Donc ^tp^tAAP = m(lambda_1\dots 0//0\lambda_r //0000)
        \item On pose \sigma_k = \sqrt{\lambda_k}. De plus, 1\leq k\leq r, le k-ime vecteur colonne de p\in R^n, q_k = \frac{1}{\sigma_k}Ap_k\in R^N
        \item monstrons que \forall k 1\leq k \leq r, ||q_k||=1, \forall k,j (q_k|q_j)=0
        \item La famille (q_1, q_r) est une famille orthonormale de R^N, mais pas forcément une base. On peut donc compléter ctte famille pour obtenir une famille avec des vecteurs supplémentaires.
        \item On note Q^n = [q_1, q_r, q_r+1, q_n]\in O(n).
        \item Montrons que AP = Q\sum
        \begin{enumerate}
            \item On calcule Q\sum = [Ap1 Apr | 0000\\] de dimension n et m-r.
            \item Calculon AP : on trouve [Ap1 Apr | Apr+1, Apn\\] de dimension n et m-r. Il faut juste montrer Apr+1=Apn = 0. Or, pr+1, pm sont des vecteurs propres de ^tAA associés à  la valeur prorpre 0. Donc pr+1,\dots pm\in ker ^tAA. Don APn = 0
        \end{enumerate}
    \end{enumerate}
\end{myproof}
\checkInfo{Méthode pour décomposer en valeurs singulières}{
\begin{enumerate}
    \item Calculer ^tAA : on s'assure que c'est symétrique
    \item Diagonaliser ^tAA en bon en rangeant les valeurs propres de ^tAA par ordre décroissant.
    \item On pose $\sigma_k = \sqrt{\lambda_k}, q_k = \frac{1}{\sigma_k}A_{pk}$
    \item Le théorème assure que la famille orthonormale complète cette famille. On complète cette famille pour obtenir une bon de R en utiisant Gram Schmitch dans Vect(q_1,\dots, q_r)^\perp.
    \item 
\end{enumerate}

}
\chapter{Espaces Hermitiens}
\section{Produit scalaire, norme}
Dans ce chapitre, on travaille sur des EV avec des complexes.
\begin{definition}
    Soit E un escpace vectoriel sur C de dimension finie ou non. C'est un produit scalaire \iff
    a) C(x+xx',y) = C(x,y) + C(x',y)
    b) C(x,y+y') pareil
    c) C(x,\lambda y) = \lambda C(x,y)
    d) C(\lambda x, y) = \bar{\lambda} C(x,y)
    e) C(y,x) = \bar{C(x,y)}
    f) c'est positif et si c'est nul, alors x est nul.
\end{definition}
b+c : y -> C(x,y) est linéaire, a+d x->C(x,y) est 1/2 linéaire : les coeff sorttent conjugués
e : C(x,x) \in R 
\begin{definition}
    Si dim E < +\infty, et C product saclaie hermitien. Alors E est un espace hermitien,  dans le cas contraire il est prohimbertien complexe,
\end{definition}
\end{document}
13/15 : 
15/16 : TP (imprimer TP)
16/17 : Maths
17/18 : EM